\title{Byte Latent Transformer: Patches Scale Better Than Tokens}

\begin{abstract}
We present the Byte Latent Transformer ({\textsc BLT}), a novel architecture for byte-level large language models that, for the first time, achieves parity with traditional tokenization-based models at scale, while offering marked enhancements in both inference speed and robustness. In {\textsc BLT}, input bytes are grouped into adaptive-length patches, which become the foundational units of computation. The boundaries of these patches are determined by assessing the entropy of the forthcoming byte, thereby dedicating greater computational effort and model resources in regions of higher data complexity. This work conducts the first flop-equated scaling analysis of byte-level architectures, evaluating models with up to 8B parameters and 4T bytes of training data. Our findings validate the viability of directly scaling byte-level models that eschew predefined vocabularies. Training and inference are rendered more efficient through the dynamic use of longer patches when the input exhibits predictability, with further benefits including qualitative gains in reasoning ability and improved generalization to rarer cases. Ultimately, at a given inference compute budget, {\textsc BLT} attains superior scaling efficiency compared to models based on fixed tokenizations, by expanding both the patch size and model capacity in tandem.
\end{abstract}

\section{Introduction}
\label{section:intro}

We present the Byte Latent Transformer~(\textbf{BLT}), a novel tokenizer-free framework designed to learn directly from raw byte sequences. Notably, BLT achieves parity with conventional tokenization-based models at scale, and further provides notable gains in both computational efficiency and model robustness (\S\ref{sec:robustness}). In contrast, current large language models (llms) are primarily trained in an end-to-end fashion, but rely on tokenization as an intermediary step. This heuristic process maps byte sequences into a predetermined vocabulary of tokens, influencing how information is encoded and introducing various limitations, such as susceptibility to domain or modality shifts~\citep{dagangetting}, heightened sensitivity to input perturbations~(\S\ref{sec:robustness}), absence of orthographic comprehension~\citep{edman2024cute}, and uneven performance across different languages~\citep{liang2023xlm,petrov2024language,limisiewicz2024myte}.

Historically, tokenization has been a necessary step, as training llms directly on raw bytes incurs substantial computational expenses, largely owing to the increased sequence lengths~\citep{xue2022byt5}. To address these challenges, earlier research has introduced more efficient self-attention mechanisms~\citep{el2020characterbert, clark2022canine} or transitioned to attention-free models~\citep{wang2024mambabyte}~(\S\ref{section:related_work}). Nonetheless, such improvements chiefly facilitate the training of \textit{small models}. When working at larger scales, the majority of computational resources in a Transformer are consumed by the sizable feed-forward network layers applied to each byte, rather than by the attention operation itself.

To optimize computational resource distribution, we introduce a dynamic, trainable approach that aggregates bytes into \textit{patches}~(\S\ref{section:patching}) and present a novel model architecture capable of integrating both byte-level and patch-level data. In contrast to tokenization, {\textsc BLT} does not rely on a predetermined set of patch tokens. Instead, it uses efficient, trainable encoder and decoder modules to project arbitrary sequences of bytes into latent patch representations. Our findings demonstrate that this strategy achieves a \textit{more} effective allocation of compute compared to approaches dependent on tokenization.

In tokenization-based LLMs, computational resources are uniformly distributed across all tokens. This approach prioritizes performance at the expense of efficiency, given that the employed compression heuristics during token induction do not consistently align with the actual prediction difficulty. The core principle underlying our architecture is the dynamic allocation of compute resources, directing more capacity to areas that demand greater modeling effort. For instance, accurately predicting the termination of a word generally does not necessitate a large transformer, since such predictions tend to involve low entropy and less complexity compared to selecting an initial word in a sentence. This guiding philosophy is embedded in the design of {\textsc BLT}~(\S\ref{section:architecture}), which incorporates three transformer blocks: two compact local models operating at the byte level and one expansive global \textit{latent transformer}~(\autoref{fig:arch_high_level}). 

To organize bytes into patches and thus assign computation adaptively, {\textsc BLT} leverages the entropy of next-byte predictions. This results in partitioning data into contextualized clusters of bytes that exhibit relatively consistent information density.

We introduce the first comprehensive flop-controlled scaling investigation of byte-level models, scaling up to 8B parameters and trained on 4T bytes, thereby demonstrating the feasibility of end-to-end large-scale training directly on bytes and removing the need for fixed-vocabulary tokenization. In terms of training efficiency under controlled computational budgets, {\textsc BLT} matches the performance\footnote{We determine model computational requirements by tallying total Floating Point OPerations (flops) performed.} of Llama 3, while at inference achieves up to 50\% reduction in flop usage~(\S\ref{section:scaling}).

Moreover, our experiments reveal that modeling inputs at the byte level leads to substantial gains, particularly for rare and uncommon inputs, enabling the model to better handle data long-tail phenomena. Compared to models relying on tokenizers, {\textsc BLT} exhibits superior robustness in the face of noisy input and demonstrates stronger character-level reasoning, as shown in tasks involving orthography, phonological phenomena, and low-resource machine translation~(\S\ref{sec:robustness}).

Additionally, the {\textsc BLT} approach facilitates simultaneous growth in both model and patch sizes under a fixed inference flop budget. Utilizing larger patch sizes not only conserves compute—allowing more capacity for the global latent transformer, which operates less frequently—but also, on average, yields significant computational savings. Our inference-flop controlled scaling studies~(\autoref{fig:fixed_inference_scaling}) display notably improved scaling patterns compared to those seen in tokenization-based model architectures.

To conclude, this work offers several key contributions: 1) We present {\textsc BLT}, a byte latent LLM framework that allocates computational resources dynamically to enhance flop efficiency; 2) We demonstrate that, at scales up to 8B, our approach matches Llama 3 in training flop usage, with the possibility to exchange marginal decreases in evaluation performance for as much as a 50\% improvement in flop efficiency; 3) {\textsc BLT} enables a novel scaling pathway for LLMs, allowing the model size to grow without increasing the fixed inference cost; 4) We provide evidence that {\textsc BLT} is more resilient to input noise and is capable of capturing sub-word information overlooked by token-level LLMs. The training and inference source code for {\textsc BLT} is publicly available at~\url{https://github.com/facebookresearch/blt}.

\section{Patching: From Individual Bytes to Groups of Bytes}
\label{section:patching}

Dividing bytes into \textit{patches} enables {\textsc BLT} to assign computational resources adaptively depending on the data’s context. As illustrated in Figure~\ref{fig:patching_types}, there are several approaches for grouping bytes into patches. More precisely, a patching function $f_p$ partitions an input byte sequence $\pmb{x}=\{x_i,|i=1,\ldots n\}$, consisting of $n$ bytes, into $m < n$ patches, denoted by $\pmb{p}=\{p_j|j=1,\ldots,m\}$, by assigning each byte $x_i$ to either 0 or 1, where a value of 1 signals the initiation of a new patch. The computational workload required by both token-based and patch-based approaches is largely a function of how many steps the main Transformer takes. Within {\textsc BLT}, this directly corresponds to the total number of patches generated by the applied patching function to encode the data. As a result, the principal determinant of computational cost for both training and inference is the average size of each patch—commonly referred to as the \textit{patch size}—for a given patching function~(\S\ref{section:flops}).

In the following sections, we present three specific patching methods: fixed-length patching, where each patch contains a constant number of bytes~(\S\ref{section:static-patch}); whitespace-based patching~(\S\ref{section:white-patch}); and a dynamic patching strategy that leverages byte-level language model entropies~(\S\ref{section:dyn-patch}). Additionally, we examine incremental patching techniques and clarify distinctions between patching and tokenization~(\S\ref{section:bpe-patch}).

\subsection{Strided Patching Every K Bytes}
\label{section:static-patch}

One of the simplest methods for organizing bytes is to divide them into patches of uniform length $k$, following the approach of MegaByte~\citep{yu2023megabyte}. This use of a fixed stride facilitates both the training and inference processes, offers a clear way to modify the mean patch size, and thus allows easy adjustment of the computational expense. Despite these advantages, this method of patch formation presents notable limitations. Primarily, computation is not adaptively distributed to the regions that require it the most: a transformer iteration $j$ could be wasted by merely predicting uninformative tokens like whitespace, or fail to allocate adequate compute to segments rich in content, such as mathematical expressions. Moreover, this strategy results in inconsistent and non-semantic partitioning of analogous byte sequences, for instance by splitting the same term differently depending on its alignment.

\subsection{Space Patching}
\label{section:white-patch}

\citet{slagle2024spacebyte} introduces a straightforward yet powerful refinement to strided patching by generating new patches whenever a space-like byte is encountered\footnote{
    Space-like bytes are defined as any byte that is not a latin character, digit, or utf-8 continuation byte. In addition, each patch must contain at least one non space-like byte. }, thereby aligning patch boundaries with linguistically meaningful units across many languages. The Space patching strategy dedicates one latent transformer operation (i.e., greater computational resources) to each word, ensuring consistent patching of words across input sequences and targeting computational effort toward predictions that are inherently more complex—often those that follow space-like bytes. For instance, the initial byte in the answer to ``Who composed the Magic Flute?\uline{\hspace{.6em}}'' poses a much greater predictive challenge compared to subsequent bytes after the initial ``M.'' This is because the first letter drastically narrows down possible continuations, whereas generating the remainder of ``Mozart'' is comparatively straightforward. Despite these advantages, Space patching falls short in handling every language or domain with flexibility and is fundamentally limited by its inability to adapt patch sizes. In what follows, we present a novel patching strategy that retains the insight that initial bytes in words are the most challenging to forecast, while also naturally enabling dynamic control over patch sizes.

\subsection{Entropy Patching: Using Next-Byte Entropies from a Small Byte LM}
\label{section:dyn-patch}

Instead of depending on heuristic methods like whitespace detection, we employ a data-driven strategy to pinpoint next-byte predictions with elevated uncertainty. We propose \textit{entropy patching}, a technique that determines patch boundaries using entropy estimates.

We train a small byte-level auto-regressive language model on the training data for {\textsc BLT} and compute next byte entropies under the LM distribution $p_{e}$ over the byte vocabulary $\mathcal{V}$:

\begin{equation}
H(x_i) = \sum_{v \in \mathcal{V}} p_{e}(x_i=v|\pmb{x}_{<i}) \log p_{e}(x_i=v|\pmb{x}_{<i})
\end{equation}

We explore two approaches for detecting patch boundaries using entropies $H(x_i)$. The initial method selects points that exceed a fixed global entropy threshold, as shown in~\autoref{fig:patching}. The alternative strategy focuses on points where the entropy is notably higher compared to the preceding value. This latter technique can be viewed as pinpointing positions that disrupt an otherwise roughly monotonically decreasing entropy trend within a patch.

\begin{align*}
    \text{Global Constraint}\hspace{1cm}H(x_t)& > \theta_g\\
    \text{Approx. Monotonic Constraint}\hspace{1cm}H(x_t)& - H(x_{t-1}) > \theta_r
\end{align*}

Patch boundaries are determined through a minimal preprocessing process performed at the dataloader stage. This approach contrasts with~\citet{nawrot-etal-2023-efficient}, where a classifier is employed to predict patch boundaries based on entropy measures. Within our experiments~(\S\ref{section:experiments}), we evaluate both approaches for differentiating low-entropy and high-entropy bytes.

\subsection{The Byte-Pair Encoding (BPE) Tokenizer and Incremental Patching}
\label{section:incremental-patching}
\label{section:bpe-patch}

A wide range of contemporary llms, such as our baseline Llama 3, employ a subword tokenization method like bpe~\citep{gage1994new, sennrich-etal-2016-neural}. In this context, we define ``tokens'' as byte sequences selected from a \textit{finite} set of possible vocabulary items established before the onset of training. In contrast, ``patches'' denote adaptively created groups of bytes without a fixed vocabulary. The primary distinction lies in the fact that, for tokens, the model is unable to directly access the raw byte-level information.

An important advancement introduced by {\textsc BLT} relative to models relying on tokenization is its reconceptualization of the balance between vocabulary size and computational expense. In conventional llms, expanding the vocabulary typically results in tokens that encapsulate more information, reducing the total number of sequential steps required for processing; however, this also expands the dimensionality of the final projection layer. This inherent trade off limits the extent to which tokenization-based models can vary token granularity and associated inference overhead. To illustrate, Llama 3 raises the mean token size from 3.7 bytes to 4.4 bytes, but this improvement necessitates a fourfold increase in the embedding table size when compared with Llama 2.

During generation, {\textsc BLT} must determine at each position in the byte sequence whether it is positioned at a patch boundary, as this choice influences the invocation of additional computation through the Latent Transformer. Importantly, this assessment is made without knowledge of subsequent, ungenerated parts of the sequence. Therefore, any patching strategy employed cannot leverage information from future bytes when determining byte sequence segmentation. Mathematically, a patching function $f_p$ exhibits the incremental patching property if it holds that:
$$f_p(\pmb{x}_{<i}) = f_p(\pmb{x})_{<i}$$
bpe does not qualify as an incremental patching scheme because the tokenization of an identical prefix may vary depending on how the sequence proceeds, thereby failing to meet the stated criterion\footnote{The inclusion of a distinct delimiter token to signal patch boundaries can make bpe compatible with incremental patching, albeit at the cost of increasing the byte sequence's length.}.

\section{{\textsc BLT} Architecture}
\label{section:architecture}
The {\textsc BLT} system consists of a major global autoregressive language model working on patch-based representations, complemented by a pair of compact local models: one responsible for converting byte sequences into patches, and the other for translating patch representations back into bytes~(\autoref{fig:arch_high_level}).

\subsection{Latent Global Transformer Model}

\textit{The Latent Global Transformer} refers to an autoregressive transformer architecture $\mathcal{G}$ composed of $l_{\mathcal{G}}$ layers. This model processes sequences of latent input patch representations, denoted $p_j$, transforming them into corresponding output patch representations, $o_j$. In this work, $j$ indexes patches, while $i$ is reserved for byte indices. The global model implements a block-causal attention masking scheme~\citep{dubey2024llama}, which confines attention to the current and preceding patches within a document. Most of the computational cost during both pre-training and inference is attributable to this model; therefore, by modulating when it is called, we can adaptively allocate computational resources across different segments of the input and output, in accordance with their respective complexities.

\subsection{Local Encoder}
\label{section:local}

\textit{The Local Encoder Model}, represented as $\mathcal{E}$, is designed as a compact transformer utilizing $l_{\mathcal{E}}$ layers, where $l_{\mathcal{E}}$ is much smaller than $l_{\mathcal{G}}$. Its principal function is to rapidly convert sequences of input bytes $b_i$ into rich patch-level representations, denoted as $p_j$. A notable innovation compared to conventional transformer architectures is the incorporation of a cross-attention layer subsequent to every transformer block. This component aggregates byte-level encodings to produce patch representations, as illustrated in~\autoref{fig:crossattn}.

To begin, each byte $b_i$ from the input is mapped through an embedding matrix of size $\mathbb{R}^{256 \times h_{\mathcal{E}}}$, resulting in embeddings $x_i$. Optionally, these can be further enhanced through the application of hash-embeddings~(\S\ref{sec:hash-emb}). 

The sequence is then processed by an alternating series of transformer layers and cross-attention layers~(\S\ref{sec:enc-cross-attn}), yielding patch representations $p_i$ that are subsequently used as input to the global transformer, $\mathcal{G}$. Each transformer block employs a \textit{local block causal} attention mechanism, where every byte is permitted to attend to a predetermined window of $w_{\mathcal{E}}$ preceding bytes; this window can generally span across moving patch boundaries but is restricted from surpassing document boundaries. The next subsections delve into a detailed discussion of both the embeddings and the mechanics of the cross-attention layer.

\subsubsection{Encoder Hash n-gram Embeddings}
\label{sec:ngram-emb}
\label{sec:hash-emb}
An essential aspect of building effective and comprehensive representations at every step $i$ is integrating context from previous bytes. In {\textsc BLT}, this is accomplished by representing the current byte $b_i$ not only as a single entity but also within the context of its surrounding n-gram. At each position $i$, we begin by generating byte-grams

\begin{align}
    g_{i,n}=\{b_{i-n + 1},\ldots, b_i\}
\end{align}

for every byte position $i$ and for $n$ ranging from three to eight.\footnote{
We do not include byte-grams of length $n$ or greater when $i<n$. }

Subsequently, we present hash $n$-gram embeddings, which assign every byte $n$-gram an index within an embedding matrix $E_{n}^{hash}$ of predetermined dimensionality, using a hash function. This procedure is conducted for each $n$ in $\{3,4,5,6,7,8\}$~\citep{bai2010learning}. The corresponding embedding is combined with the original byte embedding, then subjected to normalization prior to serving as input for the local encoder network. The process for producing the enhanced embedding is defined as follows:

\begin{align}
e_i &= x_i + \sum_{n = 3,...,8}E_{n}^{hash}(\text{Hash}(g_{i,n})) \\
\text{where, Hash}(g_{i,n}) &= \text{RollPolyHash}(g_{i,n}) \% |E_{n}^{hash}|
\end{align}

Each $e_i$ is divided by the total count of $n$-gram lengths used plus one, after which RollPolyHash, as detailed in Appendix~\ref{appendix:rollpolyhash}, is applied. Section~\ref{section:ablations} investigates how varying $n$-gram hash embedding parameters—specifically, different $n$ values and embedding table sizes—impact scaling law patterns when model flops are controlled. Besides hash-based $n$-gram embeddings, we also conducted experiments utilizing frequency-based $n$-gram embeddings, with further methodological information presented in Appendix~\ref{appendix:ngrams}.

\subsubsection{Encoder Multi-Headed Cross-Attention}
\label{sec:enc-cross-attn}
Our approach largely mirrors the input cross-attention component from the Perceiver architecture~\citep{jaegle2021perceiver}. However, rather than utilizing a predetermined set of latent vectors, we define each latent as representing an individual patch (\autoref{fig:crossattn}), ensuring that the attention is confined to only those bytes that comprise the patch itself. For each patch $p_j$, the module generates a query vector by applying a pooling operation across its corresponding byte embeddings, which is subsequently transformed using a linear projection $\mathcal{E}_{C} \in \mathbb{R}^{h_{\mathcal{E}} \times (h_{\mathcal{E}} \times U_{\mathcal{E}})}$, with $U_{\mathcal{E}}$ specifying the count of encoder cross-attention heads. Precisely, let $f_{\text{bytes}}(p_j)$ indicate the sequence of bytes associated with patch $p_j$; then we compute

\begin{align}
P_{0,j} &= \mathcal{E}_{C}(f_\text{bytes}((p_j)), f~ \text{is a pooling function} \\
P_l &= P_{l-1} + W_o\left(\text{softmax}\left(\frac{QK^T}{\sqrt{d_k}}\right)V\right) \\
\text{where } Q_j &= W_q(P_{l-1,j}), K_i = W_k(h_{l-1,i}), V_i = W_v(h_{l-1,i}) \\
h_l &= \text{Encoder-Transformer-Layer}_l(h_{l-1})
\end{align}

Here, $P \in \mathbb{R}^{n_p \times h_{\mathcal{G}}}$ denotes the representations of $n_p$ patches designated for processing by the global model. These representations are initialized by aggregating the byte embeddings $e_i$ associated with each patch $p_j$. The matrices $W_q$, $W_k$, $W_v$, and $W_o$ serve as learned projections for queries, keys, values, and outputs, respectively, with the keys and values derived by projecting byte-level features $h_i$ from the preceding layer (using $e_i$ for the initial layer). To ensure that each query $Q_j$ focuses only on information relevant to its corresponding patch $j$, we employ a specialized patch-based masking scheme where queries are restricted to attend solely to the keys and values tied to bytes within the same patch. Given the multi-head attention mechanism over $Q, K,$ and $V$, and considering that patch representations commonly have a higher dimensionality ($h_{\mathcal{G}}$) than $h_{\mathcal{E}}$, we configure $P_l$ as multiple attention heads, each of size $h_{\mathcal{E}}$, for cross-attention, then concatenate these head outputs to achieve the full $h_{\mathcal{G}}$ dimension. Furthermore, pre-LayerNorm is applied to queries, keys, and values within this cross-attention module, and no positional encodings are incorporated. Lastly, we add a residual connection around the cross-attention component.

\subsection{Local Decoder} 

In parallel to the local encoder, the local decoder $\mathcal{D}$ is designed as a lightweight transformer architecture comprising $l_{\mathcal{D}} << l_{\mathcal{G}}$ layers. Its role is to convert a sequence of global patch embeddings $o_j$ back into a series of raw bytes, $y_i$. The local decoder sequentially generates each byte, conditioning its output on previously decoded bytes and incorporating the hidden states produced by the local encoder for the sequence of bytes as input. The architecture alternates $l_{\mathcal{D}}$ layers between cross-attention and transformer blocks. Each decoding step begins with a cross-attention layer, enabling the mapping from patch representations to byte-level embeddings, followed by a transformer layer which refines these embeddings along the byte sequence.

\subsubsection{Decoder Multi-headed Cross-Attention} 

In the decoder cross-attention, the roles of the queries and key/values are interchanged i.e. the byte-representations are now the queries, and the patch representations are now the key/values. The initial byte-representations for the cross-attention are initialized as the byte embeddings from the last encoder layer i.e. $h_{l_{\mathcal{E}}}$. The subsequent byte-representations for layer $l$, $d_{l,i}$ are computed as:

\begin{align}
D_0 &= h_{l_{\mathcal{E}}} \\
B_l &= D_{l-1} + W_o\left(\text{softmax}\left(\frac{QK^T}{\sqrt{d_k}}\right)V\right), \\
  &\text{where } Q_i = W_q(d_{l-1,i}), K_i = W_k(\mathcal{D}_C(o_j)), V_i = W_v(\mathcal{D}_C(o_j)) \\
  D_l &= \text{Decoder-Transformer-layer}_l(B_l)
\end{align}

Here, as before, $W_k$ and $W_v$ are the projection matrices used to produce the key and value vectors, acting through a linear map and the split operation $\mathcal{D}_C$ over the set of final patch representations $o_j$ originating from the global model. The query projection matrix $W_q$ processes the byte-level representations $d_{l-1}$ output by the preceding decoder transformer layer (or $h_{l_{\mathcal{E}}}$ when initializing the first layer). Output vectors are then projected with $W_o$, so that $B \in \mathbb{R}^{h_{\mathcal{D}} \times n_b}$, where $n_b$ denotes the total number of output bytes. Subsequently, new decoder hidden states $D_l$ are derived by applying a decoder transformer layer to the outputs of the cross-attention module, $B$. As implemented for the local encoder’s cross-attention, our approach incorporates multi-headed attention, employs pre LayerNorm, omits positional encodings, and applies a residual connection surrounding the cross-attention layer.

\section{Experimental Setup}
\label{section:experiments}
We construct rigorously controlled experiments to evaluate {\textsc BLT} in direct comparison with models that employ tokenization, ensuring that no advantage is conferred to {\textsc BLT} through potentially increased context window length.

\subsection{Pre-training Datasets} All model variants explored in this study undergo pre-training on two primary datasets: 1) The Llama 2 dataset~\citep{touvron2023llama}, containing 2 trillion tokens amassed from diverse, publicly accessible sources, which are then subject to thorough cleaning and filtering procedures to enhance data quality; and 2) {\textsc BLT}-1T: An original dataset made up of 1 trillion tokens drawn from multiple public repositories, also incorporating a portion of the pre-training corpus distributed by Datacomp-LM~\citep{li2024datacomp}. The Llama 2 dataset is primarily utilized in scaling law analyses to identify the optimal token count, following recommendations from~\cite{dubey2024llama}, and thus inform the architectural configuration of {\textsc BLT}. In contrast, {\textsc BLT}-1T is leveraged in a full pre-training process to directly benchmark performance with Llama 3 on a range of downstream applications. Notably, neither dataset incorporates any material from Meta platforms or related services. Additionally, for baseline evaluations involving tokenizer-based models, we adopt the Llama 3 tokenizer, which possesses a 128K token vocabulary and yields better performance in our studies compared to the Llama 2 tokenizer.

\subsection{Entropy Model} For our experiments, the entropy model consists of a byte-level language model, trained on the identical data distribution as the full {\textsc BLT} architecture. Unless otherwise specified, the default configuration employs a transformer featuring 100M parameters, structured into 14 layers, with a hidden dimension size of 512, and utilizes sliding window attention spanning 512 bytes. All other hyperparameters are aligned with those used in both our local and global transformer setups. As outlined in \autoref{section:ablations}, we conducted a range of experiments varying model size, receptive field, and model architecture. Notably, in scenarios where the model's receptive field is sufficiently limited, the learned entropy model can be efficiently implemented using a lookup table.

\subsection{Entropy Threshold and Equalizing Context Length} For entropy-based patching models, we determine a patching threshold that results in the intended average \textit{patch size} over the pretraining data blend. In the case of {\textsc BLT}, in contrast to standard tokenization, the \textit{patch size} is a tunable parameter and can be set to any value, thereby directly affecting the model’s context length. To ensure a fair comparison and prevent models with larger patch sizes from benefiting due to increased context, we keep the expected total number of bytes in each batch fixed. Consequently, for configurations employing larger patch sizes, we proportionally decrease the number of patches per sequence. For Llama 2, we adopt an 8k byte context window, whereas on the {\textsc BLT}-1T dataset, the context is extended to 16k bytes per sequence on average, all while preserving an average batch size of 16M bytes.

Although the average batch size remains unchanged, dynamic patching approaches result in varying proportions of bytes to patches when assembling data batches. In our implementation of {\textsc BLT} training, we organize batches by grouping patches together, thereby minimizing the need for padding within the computationally intensive latent transformer. This approach guarantees that each batch contains an equal number of patches. To prevent memory overflow caused by exceptionally large patches during training, we pad—or, if needed, truncate—byte sequences to lengths of 12k bytes for the Llama 2 dataset and 24k bytes for the {\textsc BLT}-1T dataset.

\subsection{Entropy Model Context}
\label{section:entropy-context}
Our observations indicate that entropy patching tends to produce increasingly extensive patches in settings characterized by structured and repetitive content, such as multiple choice problems (illustrated via patching in an MMLU scenario in \autoref{fig:mmlu_patching}). These differences are attributed to reduced entropy in the repeated segments present within the entropy model context. Consequently, in the main experiment using {\textsc BLT}-Entropy with a patch size of 4.5, we refresh the entropy context using line breaks and apply an approximate monotonicity constraint, which exhibits greater resistance to “entropy drift” as context length fluctuates. Notably, this adjustment only influences the method for entropy computation, while our process for determining the entropy threshold remains unchanged.

\subsection{FLOPs Estimation} 
\label{section:flops}

Our methodology is primarily based on the transformer FLOPs calculation formulas outlined in Chinchilla~\citep{hoffmann2022training}, which include FLOPs contributions from the feed-forward layers, the query, key, value, and output (qkvo) projections within the self-attention mechanism, as well as the attention score computation and final output projection. An important deviation from prior approaches is our assumption that the input embedding layer utilizes an efficient lookup table instead of relying on dense matrix multiplication, rendering its FLOPs negligible (0-flop). In alignment with existing literature, we further approximate that the backward propagation phase consumes double the FLOPs of the forward computation.

To compute flops \textit{per byte} for {\textsc BLT} models, we add up the flops for the local encoder transformer, the global latent transformer, and the local decoder transformer, together with the cross attention blocks in the encoder and the decoder:

\begin{align}
    \text{FL}_{\text{{\textsc BLT}}} &= \text{Transf. FL}(h_{\mathcal{G}}, l_{\mathcal{G}}, m=n_{ctx}/n_p, V=0) / n_p\\
   &+ \text{Transf. FL}(h_{\mathcal{E}}, l_{\mathcal{E}}, m=w_{\mathcal{E}}, V=0) \\
   &+ \text{Transf. FL}(h_{\mathcal{D}}, l_{\mathcal{D}}, m=w_{\mathcal{D}}, V=256) \\ 
    &+ \text{Cross Attn. FL}(h_{\mathcal{E}}, l_{\mathcal{E}}, m=n_p, r=n_p/k) \cdot (k/n_p) \\ 
   &+ \text{Cross Attn. FL}(h_{\mathcal{D}}, l_{\mathcal{D}}, m=k, r=k/n_p) 
\end{align}

Here, $n_{ctx}$ refers to the length of the sequence measured in bytes, $n_p$ denotes the patch size, $r$ represents the ratio between queries and key/values, and $k$ defines the proportion between patch-dimension and byte-dimension; in other words, it is the number of individual local model partitions concatenated to create the representation for the entire model ($k=2$ in \autoref{fig:crossattn}). The parameter $V$ specifies the vocabulary size used exclusively in the output projection of the local decoder. The module may operate on either byte or patch sequences, which determines that the attention computation will utilize a different context length, $m$, accordingly. We adapt the calculation of attention-related flops for each respective model component. Detailed equations for deriving FLOPs in both the Transformer and Cross-Attention architectures are included in Appendix~\ref{section:flops-appendix}.

\subsection{Bits-Per-Byte Estimation} Perplexity only makes sense in the context of a fixed tokenizer as it is a measure of the uncertainty for each token. When comparing byte and token-level models, following previous work~\citep{xue2022byt5, yu2023megabyte, wang2024mambabyte}, we instead report Bits-Per-Byte (BPB), a tokenizer independent version of perplexity. Specifically:

\begin{align}
    \text{BPB}(x)&= \frac{\mathcal{L}_{CE}(\pmb{x})}{\ln(2)\cdot n_{\text{bytes}}}
\end{align}

Here, the cumulative cross-entropy loss reflecting the uncertainty present in the data $\pmb{x}$ is divided by both the quantity of bytes composing $\pmb{x}$ and a normalization constant.

\subsection{Transformer Architecture Hyperparameters}  
The transformer components within {\textsc BLT}, encompassing both the local and global variants, are configured to closely mirror the specifications outlined in Llama~3~\citep{dubey2024llama}. More specifically, feed-forward networks utilize the SwiGLU activation~\citep{swiglu}, rotary positional embeddings (RoPE)~\citep{RoPe} parameterized with $\theta=500000$~\citep{xiong2024effective} are applied exclusively within self-attention layers, and RMSNorm~\citep{RMSNorm} is adopted for normalization across layers. For self-attention layers employing standard, deterministic attention masks (e.g., \textit{block causal} or \textit{fixed-window block causal}), Flash attention~\citep{dao2022flashattention} is used, with a block window size of 512 when using fixed-width attention patterns. Due to the adaptive, patch-specific masking required in cross-attention modules, we leverage Flex Attention\footnote{\url{https://pytorch.org/blog/flexattention}} in order to generate optimized fused kernels, which accelerates model training considerably.

\subsection{{\textsc BLT}-Specific Hyperparameters}  
To evaluate the performance of {\textsc BLT} models, our experimental analysis proceeds along two primary lines: investigating scaling behavior and assessing downstream task performance. We examine models of varying sizes, specifically 400M, 1B, 2B, 4B, and 8B parameters. Details of the architectural hyperparameters can be found in Appendix~Table~\ref{tab:arch}.

For the first cross-attention layer in the local encoder, we initialize the queries using max-pooling. Across all {\textsc BLT} models, we utilize $500,000$ hashes with a single hash function and incorporate n-gram sizes in the range of 3 to 8. A uniform learning rate of $4e-4$ is maintained for all models. This learning rate was selected based on hyperparameter sweeps—evaluated between $1e-3$ and $1e-4$—for both token and {\textsc BLT} models at the 400M and 1B scales, which demonstrated that the same value is optimal. For scaling experiments using Llama-2 data, we adopt batch sizes in accordance with the guidelines from~\cite{dubey2024llama} or the byte-equivalent thereof. Model optimization employs the AdamW optimizer~\citep{adamw} with parameter settings $\beta_1$ = 0.9, $\beta_2$ = 0.95, and $\epsilon=10^{-8}$. We apply a linear warmup for the first 2000 steps, followed by a cosine decay of the learning rate down to zero. Further, we implement a weight decay rate of 0.1 and set a global gradient clipping threshold of 1.0.

\section{Scaling Trends}
\label{section:scaling}

We offer a comprehensive overview of the scaling behaviors exhibited by byte-level models, which provides key insights for the continued advancement of {\textsc BLT} models. Our investigation seeks to overcome several shortcomings of earlier studies on byte-level approaches in the following respects: (a) We systematically examine scaling patterns within the compute-optimal training framework, (b) We develop equivalent 8B parameter models trained on substantial datasets—extending up to 1T tokens (4T bytes)—and assess their performance on various downstream tasks, and (c) We evaluate scaling under scenarios where inference cost is held constant. Subsequent sections will focus on the distinct benefits that arise from modeling data as byte sequences.

\subsection{Parameter Matched Compute Optimal Scaling Trends}
\label{section:parameter-matched-scaling}
We conduct experiments using the Llama 2 dataset, where we train multiple \textit{compute-optimal} bpe and {\textsc BLT} models, spanning four model sizes from 1B to 8B parameters. For each model, we record the amount of training flops and assess language modeling accuracy on a representative portion of the training data mixture. The bpe models are trained following the optimal proportion of model parameters to dataset size, as prescribed by Llama~3~\citep{dubey2024llama}. This \textit{compute-optimal} approach is theoretically expected to yield peak performance on the training set given a fixed training compute~\citep{hoffmann2022training}, serving as a solid comparative standard for our analyses. Alongside each bpe model, we also train a matched {\textsc BLT} model—utilizing a Latent Transformer identical in size and architecture to its corresponding bpe Transformer—on the same dataset.

As shown in~\autoref{fig:scaling-compute-opt} (right), {\textsc BLT} models consistently perform on par with or surpass their bpe-based equivalents, with this pattern persisting as both model scale and floating point operations increase. To our knowledge, {\textsc BLT} represents the first byte-level Transformer architecture to attain comparable scaling performance to BPE-based models in compute-efficient settings. This observation thus supports our hypothesis that the ideal proportion of parameters to training compute established for bpe extends to {\textsc BLT}, or is at least approximately similar.

Both advances in architecture and the incorporation of dynamic patching are essential for aligning with bpe scaling behavior. In ~\autoref{fig:scaling-compute-opt} (left), we present a comparison of models employing space patching to Llama 3. We estimate the performance of SpaceByte \citep{slagle2024spacebyte} by using {\textsc BLT} configured for space patching, while omitting n-gram embeddings and cross-attention modules. Despite the improvements of SpaceByte relative to Megabyte, it still lags significantly behind Llama 3. The right panel of \autoref{fig:scaling-compute-opt} demonstrates the cumulative benefits derived from both architectural updates and dynamic patching techniques. When scaled, {\textsc BLT} models reach performance levels comparable to leading tokenizer-based architectures, such as Llama 3.

Additionally, our results indicate that the selection of tokenizer significantly influences model performance in tokenizer-based architectures; specifically, models utilizing the Llama-3 tokenizer demonstrate superior performance compared to counterparts trained with the Llama-2 tokenizer when evaluated on identical training datasets.

Ultimately, our {\textsc BLT} model demonstrates scaling characteristics that are intermediate between those of Llama 2 and Llama 3 when employed with substantially larger patch sizes. While the Llama 2 and Llama 3 bpe tokenizers yield average token sizes of 3.7 and 4.4 bytes respectively, {\textsc BLT} is able to replicate comparable scaling dynamics with mean patch sizes of 6 or even 8 bytes. Because inference flops decrease as the average patch size increases, utilizing a patch size of 8 bytes can result in an inference flop reduction of nearly 50\%. Moreover, as both model and dataset sizes are increased, models with greater patch sizes exhibit enhanced performance. For instance, {\textsc BLT} using an 8-byte patch size initially lags behind bpe Llama 2 at the 1B scale but surpasses it by the time the scale reaches 7B. This pattern indicates that employing such large patch sizes could yield even greater advantages at larger scales and that exploring even larger patch sizes may become increasingly viable as model parameters and computational resources expand.

\subsection{Beyond Compute Optimal Task Evaluations}
\label{section:evals}
To further explore model scaling characteristics, we continue training an 8B {\textsc BLT} model past the compute-optimal point on the {\textsc BLT}-1T dataset, which is both larger and of higher quality, and examine its performance across various widely-used classification and generation tasks. For evaluation purposes, we focus on tasks that require common sense reasoning, general world knowledge, and code generation:
\paragraph{Classification tasks} We include ARC-Easy (0-shot)~\citep{Eval_ARC}, Arc-Challenge (0-shot)~\citep{Eval_ARC}, HellaSwag (0-shot)~\citep{Eval_hellaswag}, PIQA (0-shot)~\citep{Eval_piqa}, and MMLU (5-shot)~\citep{hendrycks2020measuring}. Performance is assessed using a prompt-based likelihood method over possible answers, and we report the mean classification accuracy.
\paragraph{Coding related generation tasks:} We evaluate code generation performance by reporting pass@1 metrics on MBPP (3-shot)~\citep{Eval_MBPP} and HumanEval (0-shot)~\citep{Eval_HumanEval}, highlighting the model's effectiveness in producing Python code.

\autoref{tab:evals} presents a comparison among three models trained on the {\textsc BLT}-1T dataset: a model using the Llama 3 bpe tokenizer\footnote{The Llama 3 tokenizer is selected for its superior performance and expanded 128k vocabulary, compared to the 32k vocabulary of Llama 2.}, along with two {\textsc BLT} variants. The first variant adopts a space-patching strategy ({\textsc BLT}-Space), while the second relies on an entropy-guided patching approach ({\textsc BLT}-Entropy). The latter incorporates an approximate monotonicity constraint and refreshes its context at line breaks, in line with the procedure detailed in~\autoref{section:entropy-context}. All models were allocated an identical computational (flop) budget during training. For the {\textsc BLT}-Entropy model, we additionally reduce the inference entropy threshold from 0.6 to 0.1, an adjustment that enhances task performance albeit requiring a greater number of inference steps.

The {\textsc BLT}-Entropy model achieves superior performance compared to the Llama 3 model in 4 out of 7 tasks, despite both being trained with an equivalent quantity of bytes. This enhanced effectiveness can likely be attributed to two main factors: (1) more efficient utilization of training computation made possible by dynamic patching, and (2) explicit modeling at the byte level instead of relying on token-level representations.

Conversely, {\textsc BLT}-Space lags behind the Llama 3 tokenizer on nearly every task except one, yet it notably lowers inference flops due to its higher average patch size of 6 bytes. In contrast, both the bpe and entropy-patching based models maintain similar average patch sizes, at around 4.5 bytes on the training dataset mix. Given an equivalent training budget, the model utilizing a larger patch size processes 30\% more data than the other two, potentially moving {\textsc BLT} further from the optimal compute-efficiency threshold.

\subsection{Patches Scale Better Than Tokens} {\textsc BLT} architectures allow for the concurrent enlargement of both the model and the patch size, all while preserving a constant compute cost for training and inference and not altering the scale of the training dataset. The ability to increase patch size without constraint is a distinct characteristic of patch-based approaches, enabling them to surpass the efficiency limitations characteristic of fixed-vocabulary token-centric systems, as explored in Section~\ref{section:bpe-patch}. Utilizing larger patches reduces computational demand, making it possible to allocate those resources toward scaling up the global latent transformer, since it needs to be applied less frequently.

We perform a fixed inference scaling analysis to evaluate whether, for a given computational budget, increasing model size and using larger patches with fewer processing steps yields superior performance compared to employing smaller models with more steps. Using models with 400 million and 3.6 billion parameters initialized with the Llama 2 tokenizer, we identify flop-matched models using the Llama 3 tokenizer as well as {\textsc BLT}-Entropy architectures featuring mean patch sizes of 6 and 8 bytes over the mixed training data (refer to~\autoref{tab:fixed-inf-params} for specific configurations). For those models utilizing patch size 8, we employ 3 encoder layers rather than just 1. Each variant is trained under multiple training flop constraints.

As illustrated in \autoref{fig:fixed_inference_scaling}, {\textsc BLT} models demonstrate superior scaling behavior compared to architectures that rely on tokenization, across both categories of inference flop. While BPE models tend to outperform at limited training compute, {\textsc BLT} models quickly overtake them as training extends just beyond the regime optimized for compute efficiency. In many scenarios, allocating greater resources to the pretraining stage—despite it being a one-off expenditure—can yield models with substantially improved performance under a fixed inference compute constraint. This approach is exemplified by the 8B model class, such as Llama 3.1, which was trained on data volumes two orders of magnitude larger than the compute-optimal quantity for its size.

The point at which {\textsc BLT} begins to outperform token-based approaches has moved slightly nearer to the compute-optimal threshold as the flop class increases; specifically, the shift is from three times to 2.5 times the compute-optimal expenditure. Likewise, for models with patch size 8, we observe that in higher flop classes, the scaling curve becomes steeper and these models surpass their counterparts more quickly. As elaborated in Section~\ref{section:parameter-matched-scaling}, increasing the patch size appears to align the performance of these models with BPE models as the overall model scale grows. This phenomenon can be partially explained by the diminishing portion of total flops consumed by the byte-level Encoder and Decoder, which seem to scale at a slower rate than the Latent Transformer component. For instance, scaling model parameters twenty-fold from 400M to 8B approximately results in just a two-fold increase in the local parameters for {\textsc BLT}. This distinction is significant because patch size only impacts the flops associated with the patch Latent Transformer, leaving the byte-level modules unaffected. For this reason, the ratio of {\textsc BLT}-Entropy ps=8 model size to that of Llama 2 increased from 1.6x to 1.7x upon transitioning to the larger model configuration.

In conclusion, our analysis of patch-length scaling reveals that the {\textsc BLT} patch-based architecture benefits from joint increases in patch size and model size, yielding superior scaling behavior. Notably, these scaling improvements not only hold but may even become more pronounced as models grow larger.

\section{Byte Modeling Improves Robustness}
We further assess how {\textsc BLT}'s robustness compares to that of token-based models, which do not have intrinsic byte-level information, and introduce a methodology for converting pretrained token-based models into byte-oriented ones.
\label{sec:robustness}

\subsection{Character-Level Tasks}

Initial interest in developing byte-level models stemmed from their capacity to handle noise at the byte level in input data and their intrinsic understanding of token subcomponents—areas where models reliant on tokenizers tend to be less effective. To systematically assess these strengths, we conduct supplementary tests on benchmarks designed to quantify not only resilience to noisy inputs but also recognition and processing of character-level features, encompassing English, multiple languages, numerical digits, and phonetic symbols. The corresponding findings are reported in Table~\ref{tab:char_tasks}.

\paragraph{Noisy Data} To assess the resilience of {\textsc BLT} in comparison with tokenizer-based approaches, we generate perturbed versions of the benchmark classification datasets discussed in Section~\ref{section:evals}. We introduce text variation using five unique character-level noise techniques: (a) \textit{AntSpeak}: Converts the entire text into uppercase letters separated by spaces. (b) \textit{Drop}: Randomly deletes 10\% of the characters. (c) \textit{RandomCase}: Randomly assigns half the characters to uppercase and the other half to lowercase. (d) \textit{Repeat}: For 20\% of the characters, increases their frequency by repeating them up to four times. (e) \textit{UpperCase}: Modifies the text so all characters appear in uppercase. Each perturbation is separately applied to the prompt, completion, or both during the evaluation phase, and results are averaged for reporting. As shown in Table~\ref{tab:char_tasks}, we present our findings for the noised HellaSwag~\citep{Eval_hellaswag} dataset. Our analysis reveals that {\textsc BLT} consistently demonstrates superior robustness compared to tokenizer-based baselines, achieving on average an 8-point margin over an equivalently-trained model and surpassing the Llama 3.1 model even though the latter is trained on significantly larger data.

\paragraph{Phonology - Grapheme-to-Phoneme (G2P)} We evaluate how effectively {\textsc BLT} converts sequences of graphemes—that is, characters encoding a word—into corresponding phonemic transcriptions that represent the word’s pronunciation. Table \ref{tab:char_tasks} displays the outcomes of this G2P evaluation conducted in a 5-shot format using Phonology Bench~\citep{suvarna2024phonologybench}. Our analysis shows that {\textsc BLT} achieves superior results compared to the baseline Llama 3 1T model, which employs a tokenizer-based approach for this task.

\paragraph{CUTE}
To evaluate the character-level capabilities of {\textsc BLT}, we test it on the CUTE benchmark~\citep{edman2024cute}, which includes a range of tasks that fall into three main groups: tasks involving compositional understanding, those testing orthographic similarity, and sequence manipulation tasks. The CUTE benchmark is particularly difficult for most tokenizer-based language models, which typically know the spelling of their internal tokens but find it challenging to leverage that knowledge for actual text manipulation. As presented in Table~\ref{tab:char_tasks}, {\textsc BLT}-Entropy exceeds the performance of both BPE Llama 3 models by over 25 points across these tasks. Notably, our model is particularly strong at manipulating character sequences, obtaining near-perfect scores (99.9\%) on spelling-related tasks. These pronounced gains—despite {\textsc BLT} being trained on only 1/16th the amount of data used for Llama 3.1—suggest that BPE models have significant difficulty in acquiring fine-grained character-level representations. Figure \ref{fig:cute_examples} highlights examples where the Llama 3 tokenizer model performs suboptimally, whereas {\textsc BLT} is much more successful. The only areas where BPE models outperform are word deletion and insertion, tasks for which token-based architectures are well suited. However, the margin remains small, and character-based models like ours may bridge this gap more easily than BPE models could reverse the process. For all tasks, we adopt a consistent evaluation protocol, utilizing the original Huggingface prompts. It is possible that BPE models could see improved results with further prompt engineering.

\paragraph{Low Resource Machine Translation} We assess {\textsc BLT} for translation tasks involving both directions across six prominent language families and twenty-one low-resource languages, spanning multiple scripts, utilizing the FLORES-101 benchmark~\citep{goyal2022flores}. Table~\ref{tab:flores} provides the SentencePiece BLEU scores for these evaluations. The findings reveal that {\textsc BLT} surpasses a counterpart trained with the Llama 3 tokenizer, exhibiting an overall gain of 2 BLEU points for translations into English and a 0.5-point improvement for translations from English. For widely used language pairs, {\textsc BLT} delivers results that are on par with or marginally superior to Llama 3. Importantly, {\textsc BLT} consistently yields better performance than Llama 3 on several low-resource language pairs, highlighting the robustness of byte-based modeling in capturing and generalizing across underrepresented byte patterns.

\subsection{Training {\textsc BLT} from Llama 3}
\label{sec:distillation}
We investigate a methodology where {\textsc BLT} models benefit from pre-existing, pre-trained tokenizer-based architectures to facilitate accelerated and improved training convergence. This is accomplished by initializing {\textsc BLT}'s global transformer parameters with those extracted from a pre-trained Llama 3.1. Following initialization, the global transformer’s weights are finetuned using a learning rate set at one-tenth of that applied to both the local encoder and local decoder components, over the optimal training steps for Llama 3. Table~\ref{tab:distill} contrasts this strategy against a baseline {\textsc BLT}. Results clearly indicate that initializing {\textsc BLT} from Llama 3.1 yields substantial performance improvements over both standard Llama 3 and the baseline {\textsc BLT}, with all models trained for an equal number of flops. Furthermore, in comparison to the {\textsc BLT}-Entropy variant—see Table~\ref{tab:evals}—which involved pre-training on a much larger dataset comprising 1T tokens, the Llama 3.1-initialized {\textsc BLT} demonstrates even better results on the MMLU benchmark. This indicates that such an approach can offer meaningful reductions in training computational cost.

This approach can be interpreted as converting models that rely on tokenizers into ones that do not, thus enabling the transformation of a pre-trained LLaMA 3.1 model into a {\textsc BLT} model. For a thorough comparison, we present results for the original LLaMA 3.1 model, which was trained on 15T tokens, in Table \ref{tab:distill}, and benchmark it alongside the LLaMA-derived {\textsc BLT}. While our model only displays a minor decrease in performance on MMLU and HumanEval, it experiences a more substantial reduction across several other benchmarks. These findings highlight the necessity for further investigation, particularly to refine data mixtures and tune hyperparameters to make more effective use of the pre-trained model’s capabilities.

\section{Ablations and Discussion}
\label{section:ablations}
Here, we present a series of ablation studies to validate our design decisions regarding the {\textsc BLT} architecture, as well as the patching methodology and selection of hyper-parameters used for training the {\textsc BLT} model with 8B parameters on the {\textsc BLT}-1T dataset.

\paragraph{Entropy Model Hyper-parameters} To investigate how the size of the entropy model and the length of its context window impact scaling behavior, we experiment with byte-level entropy transformer architectures that range in size from 1 million to 100 million parameters and employ context windows of lengths between 64 and 512. In \autoref{fig:entropy_ablation}, we display curves relating bpb to training flop counts, derived from our $400m$ and $1b$ {\textsc BLT} models trained using the Llama-2 dataset. The results reveal that increasing either model size or context window enhances scaling outcomes, although improvements become marginal beyond the 50 million parameter scale.

\paragraph{Types of Patching}

We conduct an ablation study on the four patching strategies described in Section~\ref{section:patching}, namely: 1) Strided Patching using strides of 4 and 6, 2) Whitespace-based Patching, 3) Patching guided by BPE Tokenizer segmentation as per the Llama 3 tokenizer, and 4) Entropy-driven patching implemented with a lightweight byte-level language model.

Although dynamic patching decreases the actual sequence length, we regulate sequence length to ensure that every patching approach has an equivalent context length. Consequently, each model is exposed to an identical expected number of bytes per sequence during both training and inference, eliminating potential confounds related to context size differences. As illustrated in Figure~\ref{fig:scaling-compute-opt}, these ablation results reveal that all alternative patching strategies surpass static patching in performance. Notably, space patching performs nearly as well as dynamic entropy based patching.

In \autoref{tab:evals_ablation}, we report benchmark results for {\textsc BLT} models, directly comparing approaches including tokenizer-based modeling, space patching, and entropy-based patching, all trained on the Llama 2 dataset for an appropriate number of iterations~\citep{dubey2024llama}. While space patching constitutes a more straightforward method that does not require the entropy model to be executed dynamically during training, our analysis demonstrates that the improvements associated with entropy-based patching, as noted in the scaling experiments (see Section~\ref{section:scaling}), persist in downstream benchmark assessments.\footnote{Space patching outcomes are derived from preliminary runs lacking cross-attention; however, comparable patterns emerge in settings where cross-attention is included.}

\paragraph{Cross-Attention} 
In~\autoref{tab:crossattn_ablations_1b}, we systematically remove cross-attention at different stages within both the encoder and decoder components of {\textsc BLT} to assess its impact. For the encoder, we evaluate three methods for initializing the cross-attention queries: 1) using a single shared learned embedding across all global states, 2) employing a hash embedding based on the patch's byte sequence, and 3) deriving the pooled encoder hidden states of the patch bytes at a specific encoder layer.

Our results indicate that incorporating cross-attention within the \textit{decoder} yields the best performance. When implemented in the encoder, cross-attention offers only marginal gains, and these occur primarily when query initialization employs pooling. Moreover, cross-attention proves especially beneficial on Common-Crawl data, with its positive impact becoming more pronounced as patch size increases.

\paragraph{n-gram Hash Embeddings} 

We experiment with various n-gram hash embedding vocabulary sizes—specifically 0, 100K, 200K, and 400K—and summarize the outcomes in Table~\ref{tab:ngram_ablations_1b}. Our analysis demonstrates that hash embeddings consistently yield performance benefits across all evaluated domains, with especially pronounced improvements observed on Wikipedia and Github datasets (a 0.04 bpb improvement versus 0.01 bpb after 15k steps at the 8B parameter scale). At the 8B model size, reducing the number of hashes from 500K to 300K results in only a marginal change of approximately 0.001 bpb at 15k training steps. These findings suggest that hash embeddings are crucial for {\textsc BLT} to reach parity with tokenizer-based models, although performance improvements tend to plateau once 300K hash entries are reached. Moreover, the observed advantages provided by hash embeddings appear largely independent of those produced by cross-attention, as the two contribute distinct gains across different dataset types.

\paragraph{Local Model Hyperparameters} Table \ref{tab:local_layers} presents an ablation study examining different configurations of the number of layers in both the local encoder and decoder. In conjunction with hash n-gram embeddings, {\textsc BLT} achieves strong performance even when utilizing an extremely minimal encoder—consisting of just a single layer—while benefiting from a decoder of greater depth.

\section{Related Work}
\label{section:related_work}

\textbf{Character-Level RNNs:} The task of Character Language Modeling has maintained its significance since the advent of early neural architectures~\citep{sutskever2011generating,mikolov2012subword,graves2013generating}, largely due to their intrinsic ability to natively handle out-of-vocabulary terms, eliminating the necessity for explicit back-off strategies. For instance, \cite{kim2016character} proposed a system utilizing convolutional and highway networks to process input characters, which are subsequently passed to LSTM-based RNNs. This approach achieved comparable performance to leading RNN-based language models at the time on English datasets, and even surpassed them on morphologically complex languages—a notable benefit of character-level large language models. Expanding on these findings, \cite{kenter2018byte} demonstrated that byte-level LSTM networks were superior to traditional word-level models for machine comprehension tasks in highly inflected languages such as Turkish and Russian. Similarly, \cite{zhang2015character} developed character-level convolutional models for classification applications, noting their improved performance over word-based approaches for select tasks. Building on the idea of deeper linguistic structure, \citet{chung2019hierarchical} incorporated hierarchical LSTM architectures featuring boundary-detector mechanisms at each stage to infer implicit hierarchical structure within text, thereby further enhancing the performance on character-level language modeling. Alternatively, ByteNet by \cite{kalchbrenner2016neural} implements convolutional neural layers over character sequences in place of attention mechanisms to facilitate machine translation.

\textbf{Character-Level Transformers:} The introduction of attention mechanisms in transformer architectures~\citep{vaswani2017attention}, along with advancements in subword tokenization~\citep{sennrich-etal-2016-neural}, resulted in substantial gains for neural language models and established benchmarks. Nonetheless, choosing to model data at the word or subword level imposes an inherent inductive bias regarding the abstraction scale at which the model operates. In pursuit of integrating the advantages of transformer-based models with earlier encouraging outcomes in character-level language modeling, \citet{al2019character} explored stacking deep transformer networks and incorporating auxiliary losses, achieving better results compared to previous LSTM-based character language models. Despite these improvements, their character-level transformers lagged notably behind large language models trained on word-level input. Additionally, GPT-2~\citep{radfordgpt2} reported that for extensive datasets, such as the 1 Billion Word Benchmark, models operating at the byte level fell short of matching the performance of their word-level counterparts.

As shown in \cite{Choe2019BridgingTG}, transformer-based language models at the byte level can surpass the performance of subword-based models with a similar number of parameters, but at the cost of significantly higher computational requirements and extended training times. Likewise, \cite{el2020characterbert} introduce a BERT variant named CharFormer, which constructs word-level representations via convolutional operations on character embeddings. While this approach yields gains within the medical text domain, it also entails substantially increased computational resources. The CANINE model \cite{clark2022canine}, comprising 150M parameters in an encoder-only configuration, processes text as raw character sequences. Its architecture integrates a deep transformer—serving a purpose analogous to our global model—and employs both a local transformer and strided convolutions for character-level downsampling. As a result, CANINE outperforms similarly-sized token-based encoder-only models (e.g., mBERT) in various multilingual benchmarks. ByT5~\citep{xue2022byt5} investigates byte-level encoder-decoder models that eschew patching mechanisms altogether. Although this configuration shows enhanced robustness to input noise and remains competitive with tokenized models, requiring only a quarter of the data, the absence of patching forces the models to perform intensive attention calculations over each byte—dramatically increasing computational burden. The challenge with byte-level modeling arises from the greater input sequence length, which hampers the efficiency and scalability of such models. More recently, work utilizing the Mamba Architecture~\citep{gu2023mamba}, known for maintaining constant-size memory across extended contexts, demonstrates that a patchless, byte-level Mamba-based model can outperform transformer-based byte-level models in flop-equivalent settings at the 350M parameter scale, particularly when measured by bits-per-byte across multiple datasets~\cite{wang2024mambabyte}.

\textbf{Patching-based approaches:} By utilizing patching techniques, it is possible to significantly reduce the large number of flops that are typical of byte-level LLMs, without sacrificing performance. Various studies have reported promising results on models with modest sizes and limited training data. \cite{nawrot-etal-2022-hierarchical} introduced static patching with downsampling and upsampling mechanisms, proposing the hourglass transformer, which achieved superior performance compared to other byte-level models at the 150M parameter level. Building upon this, \cite{nawrot-etal-2023-efficient} introduced more sophisticated dynamic patching strategies, such as a learnable boundary-predictor trained end-to-end, a boundary-predictor with guidance from specific tokenizers, and an entropy-driven patching strategy resembling {\textsc BLT}. Their experiments demonstrated that these methods could outperform standard transformer models of the time on language modeling benchmarks with 40M parameters and around 400M tokens. Additionally, \cite{lester2024training} explored models trained on sequences encoded through arithmetic coding, achieving compression efficiencies beyond those offered by BPE. By implementing the equal-info windows method, they reported improvements over byte-level baselines in language modeling tasks, although their results did not surpass those of subword-based baselines.

Our research is primarily inspired by and most similar to MegaByte~\citep{yu2023megabyte}, a decoder-only causal LLM that employs a predetermined fixed patching strategy and concatenation of representations to transform bytes into patches, subsequently applying a local model within the decoder to revert patches to bytes. Their results indicate that MegaByte achieves parity with tokenizer-based models at the 1B parameter scale using a corpus of approximately 400B bytes. Throughout our experimental evaluation, we include ablation studies of MegaByte and observe that the static patching approach falls short of the state-of-the-art compute-efficient tokenizer-based models under fixed FLOP budgets, and we illustrate how {\textsc BLT} addresses this performance gap. \cite{slagle2024spacebyte} similarly note MegaByte’s limitations and propose an enhanced static patching methodology by targeting whitespace and related space-like bytes, in addition to introducing a local encoder model. Their work demonstrates performance gains over tokenizer-based transformer models under comparable compute constraints in particular domains such as Github and arXiv at the 1B scale. Our experiments with their approach corroborate these findings but also highlight the necessity for continued architectural innovation to further advance byte-level models in order to effectively rival the leading token-based frameworks like Llama 3.

\section{Limitations and Future Work}

In this study, our architectural decisions are informed by training models for the number of steps considered optimal for Llama~3~\citep{dubey2024llama}. Notably, these scaling laws were derived with BPE-level transformers in mind, which could result in suboptimal relationships between data and parameter sizes when applied to {\textsc BLT}. Determining scaling laws tailored specifically for {\textsc BLT} remains an avenue for future investigation, and could reveal more advantageous scaling behavior unique to our approach. Furthermore, many of our experiments have been limited to models with up to 1B parameters; it is conceivable that the optimal design choices may differ when extending to 8B parameters or larger, thereby potentially enabling further performance gains at greater model sizes.

Current transformer libraries and codebases are primarily optimized for high efficiency with tokenizer-based transformer models. Although we conduct theoretical flop-matched comparisons and leverage some optimized solutions—such as FlexAttention—to support layers that differ from standard transformer designs, our implementations might still lag behind tokenizer-based counterparts in actual wall-clock performance and could be further accelerated through additional optimizations.

Although {\textsc BLT} currently relies on an independently trained entropy model for patching, integrating the learning of the patching model into an end-to-end training process represents a promising avenue for subsequent research. In Section \ref{sec:distillation}, we report preliminary results which suggest that it is feasible to “byte-ify” tokenizer-based models—such as Llama 3, which have been trained on datasets exceeding 10T tokens—by initializing and keeping the global transformer's weights fixed. Pursuing this line of investigation could lead to techniques that preserve the strengths of bytefying, while also achieving superior performance compared to these tokenizer-based models, all without necessitating training from the ground up.

\section{Conclusion}
\label{section:conclusion}

This work introduces the Byte Latent Transformer (\textbf{BLT}), an innovative framework that challenges the standard reliance on static-vocabulary tokenization in large language models. Through the implementation of a flexible, learnable mechanism for aggregating bytes into patches, {\textsc BLT} enables adaptive computational allocation that reflects the underlying complexity of the data. This approach results in notable enhancements in both computational efficiency and model robustness. Comprehensive scaling experiments reveal that {\textsc BLT} achieves competitive performance with tokenization-based architectures such as Llama 3 at sizes up to 8B parameters and across 4T bytes of training data, while offering the possibility to decrease inference flops by as much as 50\% with minimal impacts on key evaluation benchmarks. Additionally, {\textsc BLT} introduces a novel axis for architectural scaling, permitting parallel enlargement of both the model size and patch dimensions within a static inference cost constraint. This capability is especially beneficial for typical compute environments. By operating directly on byte-level input, {\textsc BLT} enhances the model’s proficiency in processing less frequent and noisier data, resulting in superior resilience to non-standard or corrupted inputs and richer representations of sub-word elements. Collectively, these findings identify {\textsc BLT} as a compelling substitute to conventional tokenization-centric models, supporting the development of language models that are not only more efficient but also more robust and scalable.

\end{document}